\documentclass[11pt,a4paper]{article}

\usepackage[utf8]{inputenc}
\usepackage[T1]{fontenc}
\usepackage[margin=2.5cm]{geometry}
\usepackage{hyperref}
\usepackage{amsmath}
\usepackage{amssymb}
\usepackage{booktabs}
\usepackage{longtable}
\usepackage{array}
\usepackage{xcolor}
\usepackage{listings}
\usepackage{graphicx}
\usepackage{enumitem}
\usepackage{titlesec}
\usepackage{fancyhdr}
\usepackage{microtype}
\usepackage{parskip}

% Colours
\definecolor{ocha-blue}{RGB}{0,114,188}
\definecolor{code-bg}{RGB}{245,245,245}
\definecolor{code-comment}{RGB}{100,100,100}
\definecolor{critical-red}{RGB}{185,28,28}
\definecolor{high-orange}{RGB}{180,83,9}

% Hyperref setup
\hypersetup{
  colorlinks=true,
  linkcolor=ocha-blue,
  urlcolor=ocha-blue,
  citecolor=ocha-blue,
  pdftitle={Humanitarian Crisis Neglect Index - Technical Documentation},
  pdfauthor={Alles GUT},
}

% Code listings
\lstset{
  backgroundcolor=\color{code-bg},
  basicstyle=\ttfamily\small,
  breaklines=true,
  frame=single,
  rulecolor=\color{gray!50},
  commentstyle=\color{code-comment},
  keywordstyle=\color{ocha-blue}\bfseries,
  stringstyle=\color{critical-red},
  numbers=left,
  numberstyle=\tiny\color{gray},
  numbersep=6pt,
  xleftmargin=12pt,
  tabsize=4,
  showstringspaces=false,
}

\titleformat{\section}{\large\bfseries\color{ocha-blue}}{}{0em}{\thesection\quad}
\titleformat{\subsection}{\normalsize\bfseries}{}{0em}{\thesubsection\quad}
\titleformat{\subsubsection}{\normalsize\itshape\bfseries}{}{0em}{\thesubsubsection\quad}

\pagestyle{fancy}
\fancyhf{}
\lhead{\small\textcolor{ocha-blue}{\textbf{Humanitarian Crisis Neglect Index}}}
\rhead{\small Technical Documentation}
\cfoot{\thepage}
\renewcommand{\headrulewidth}{0.4pt}

\newcommand{\code}[1]{\texttt{#1}}
\newcommand{\apiendpoint}[1]{\texttt{\textcolor{ocha-blue}{#1}}}
\newcommand{\param}[1]{\texttt{\textit{#1}}}


\begin{document}

\begin{titlepage}
  \centering
  \vspace*{3cm}
  {\Huge\bfseries\color{ocha-blue} Humanitarian Crisis Neglect Index \par}
  \vspace{0.5cm}
  {\Large Technical Documentation \par}
  \vspace{0.8cm}
  \rule{10cm}{0.5pt}\par
  \vspace{0.8cm}
  {\large\textit{Identifying Severe but Underfunded Humanitarian Crises} \par}
  \vspace{2cm}
  {\normalsize
    \begin{tabular}{rl}
      \textbf{Team}  & Alles GUT \\[4pt]
      \textbf{Task} & Geo-Insight: Which Crises Are Most Overlooked? \\[4pt]
      \textbf{Date}    & April 2026 \\[4pt]
    \end{tabular}
  \par}
  \vfill
\end{titlepage}

% --- Table of contents -------------------------------------------------------
\tableofcontents
\newpage

% -----------------------------------------------------------------------------
\section{Introduction}
% -----------------------------------------------------------------------------

\subsection{Motivation}

Not every humanitarian crisis receives funding proportional to its documented need.
Some situations are severe, chronic, and well-documented - yet year after year they
receive far less than their requested requirements.
Others may be similarly dire but are acute or newly escalating, making them
``structurally'' overlooked rather than recently neglected.

The \textbf{Humanitarian Crisis Neglect Index (HCNI)} is a data-driven platform that
surfaces these overlooked crises by combining two complementary signals:

\begin{itemize}[nosep]
  \item \textbf{Crisis severity} --- the scale and urgency of documented humanitarian need,
        derived from people-in-need counts, food-insecurity phase data, and conflict events.
  \item \textbf{Funding gap} --- how far the current funding coverage falls short of
        requirements, computed from OCHA's Financial Tracking Service (FTS).
\end{itemize}

A crisis scores high on the neglect index when it is \emph{simultaneously} severe and
underfunded relative to its cluster peers.
The system is designed for decision support, not automated allocation: outputs are
ranked and explained so that a humanitarian analyst or donor advisor can interrogate,
adjust, and act on them.

\subsection{Design Principles}

\begin{enumerate}[nosep]
  \item \textbf{Transparency} --- every component of the score is visible and adjustable.
  \item \textbf{Uncertainty awareness} --- bootstrap confidence intervals surface ranking
        instability; missing data is flagged, not silently imputed.
  \item \textbf{Within-cluster fairness} --- all percentile ranks compare crises only
        against peers in the same humanitarian cluster (e.g.\ Health vs.\ Shelter).
  \item \textbf{Structural vs.\ acute neglect} --- a six-year FTS history distinguishes
        chronic underfunding from a recent deterioration.
  \item \textbf{Human-in-the-loop} --- the Claude-powered chatbot acts as an analytical
        co-pilot; it never makes allocation decisions autonomously.
\end{enumerate}

% -----------------------------------------------------------------------------
\section{Methodology}
% -----------------------------------------------------------------------------

\subsection{Neglect Index Formula}

For each crisis $c$ (identified by country, humanitarian cluster, and year), the
neglect index $\mathcal{N}(c)$ is defined as:

\begin{equation}\label{eq:neglect}
  \mathcal{N}(c) \;=\; \widetilde{S}_w \cdot \text{severity}(c)
                  \;+\; \widetilde{G}_w \cdot \bigl(1 - \text{coverage\_rank}(c)\bigr)
\end{equation}

\noindent where $\widetilde{S}_w$ and $\widetilde{G}_w$ are the severity weight and
funding-gap weight normalised to sum to 1:
\[
  \widetilde{S}_w = \frac{S_w}{S_w + G_w}, \qquad
  \widetilde{G}_w = \frac{G_w}{S_w + G_w}.
\]
Default weights: $S_w = 0.6$, $G_w = 0.4$.

\subsection{Severity Score}

The severity term itself is a weighted combination of up to three signals.
Exactly which signals are used depends on data availability (Cases A--D):

\begin{description}[labelindent=0.5cm, labelsep=1em]
  \item[Case A (IPC + conflict events available)]
    \[
      \text{severity}(c) \;=\; \widetilde{N}_w \cdot \text{need\_rank}(c)
                          \;+\; \widetilde{I}_w \cdot \text{ipc\_severity}(c)
                          \;+\; \widetilde{E}_w \cdot \text{events\_rank}(c)
    \]
  \item[Case B (IPC only)]
    \[
      \text{severity}(c) \;=\; \widetilde{N}_w \cdot \text{need\_rank}(c)
                          \;+\; \widetilde{I}_w \cdot \text{ipc\_severity}(c)
    \]
  \item[Case C (conflict events only)]
    \[
      \text{severity}(c) \;=\; \widetilde{N}_w \cdot \text{need\_rank}(c)
                          \;+\; \widetilde{E}_w \cdot \text{events\_rank}(c)
    \]
  \item[Case D (neither IPC nor events)]
    \[
      \text{severity}(c) \;=\; \text{need\_rank}(c)
    \]
\end{description}

In all cases the sub-weights $\widetilde{N}_w$, $\widetilde{I}_w$, $\widetilde{E}_w$
are re-normalised to sum to 1 using the available components.
Default sub-weights: $N_w = 0.5$, $I_w = 0.4$, $E_w = 0.1$.

\subsection{Component Definitions}

\subsubsection{need\_rank}

Percentile rank of the number of \emph{people in need} for crisis $c$ among all
crises in the same humanitarian cluster.
Ranks are computed on the aggregate dataset after collapsing multi-year rows into a
single (country, cluster) observation.

\subsubsection{ipc\_severity}

The Integrated Food Security Phase Classification (IPC) records the population share
in each phase 1--5 for the ``Food Security'' cluster.
The IPC severity score maps the average phase $\bar{p}$ to a $[0, 1]$ scale with
exponential weighting that penalises famine (phase 5) strongly:

\begin{equation}\label{eq:ipc}
  \text{ipc\_severity} \;=\; \frac{2^{\,\bar{p} - 1} - 1}{15}
\end{equation}

\noindent This gives $\text{ipc\_severity} = 0$ at phase 1 (minimal) and
$\text{ipc\_severity} = 1$ at phase 5 (famine).

\subsubsection{events\_rank}

Percentile rank of the number of ACLED conflict events targeting civilians in the
country and year of crisis $c$, computed within the same humanitarian cluster.

\subsubsection{coverage\_rank}

Percentile rank of the funding coverage ratio:
\[
  \text{coverage}(c) \;=\; \min\!\left(\frac{\text{funding}(c)}{\text{requirements}(c)},\; 1\right)
\]
A high \code{coverage\_rank} (coverage better than peers) produces a low
contribution to the neglect index via the $(1 - \text{coverage\_rank})$ term.

\subsection{Structural Neglect Signals}

Using six years of FTS history (2020--2025), each (country, cluster) pair is
characterised by four structural signals:

\begin{description}[labelindent=0.5cm, labelsep=1em]
  \item[\texttt{consecutive\_years\_underfunded}]
    The length of the current streak of years with coverage $<50\%$.
  \item[\texttt{structural\_neglect\_score}]
    \[
      \text{sn\_score} \;=\; 0.6 \cdot \frac{\min(\text{streak}, 5)}{5}
                       \;+\; 0.4 \cdot \text{pct\_years\_underfunded}
    \]
  \item[\texttt{coverage\_trend}]
    Linear regression slope of coverage over the six-year window.
    A negative value indicates worsening coverage.
  \item[\texttt{neglect\_type}]
    One of: \textbf{structural} ($\geq$3 consecutive years + $\geq$60\% of window
    underfunded), \textbf{worsening} (negative trend + currently underfunded),
    \textbf{acute}, \textbf{improving}, or \textbf{adequate}.
\end{description}

\subsection{Bootstrap Uncertainty Estimation}

When \param{n\_bootstrap} $> 0$ (default 200), the scoring procedure is repeated on
random sub-samples of the crisis pool.
The resulting rank distribution is summarised as a 90\% confidence interval
$[\text{rank}_{5\%},\, \text{rank}_{95\%}]$.
A wide interval signals that a crisis's rank is sensitive to which other crises happen
to be in the dataset---a meaningful caveat for sparse clusters.

\subsection{Priority Labels}

Crises are labelled by neglect index thresholds (configurable):

\begin{center}
  \begin{tabular}{lll}
    \toprule
    Label & Default threshold & Interpretation \\
    \midrule
    \textcolor{critical-red}{\textbf{Critical}} & $\mathcal{N} \geq 0.80$ & Immediate attention warranted \\
    \textcolor{high-orange}{\textbf{High}}     & $0.60 \leq \mathcal{N} < 0.80$ & Elevated concern \\
    \textbf{Medium}  & $0.40 \leq \mathcal{N} < 0.60$ & Monitor closely \\
    \textbf{Low}     & $\mathcal{N} < 0.40$ & Relatively well covered \\
    \bottomrule
  \end{tabular}
\end{center}

% -----------------------------------------------------------------------------
\section{System Architecture}
% -----------------------------------------------------------------------------

The system follows a three-tier architecture: a data layer of static CSV/XLSX files
plus a live HDX HAPI feed, a Python FastAPI backend, and a React frontend.

\begin{figure}[h]
\centering
\begin{minipage}{0.9\textwidth}
\begin{lstlisting}[language={},numbers=none,frame=single,
  caption={High-level architecture overview}]
+----------------------------------------------------------+
|                React Frontend  (:3000)                   |
|  CsvViewer . TsneChart . Chat/Weights . WorldMap         |
|                  CounterfactualSlider                    |
+----------------------+-----------------------------------+
                       |  HTTP / JSON  (:8000)
+----------------------v-----------------------------------+
|              FastAPI Backend  (:8000)                    |
|   main.py  -  chatbot.py  -  scorer.py  -  hdx_client   |
+----------------------+-----------------------------------+
        +--------------+--------------+
        |                             |
   CSV / XLSX files             HDX HAPI (live)
   (data/ directory)            api.hpc.tools
\end{lstlisting}
\end{minipage}
\end{figure}

\subsection{Backend (\texttt{api/})}

\begin{description}[labelindent=0.5cm, labelsep=1em]
  \item[\texttt{main.py}] FastAPI application.  Defines all HTTP endpoints,
    loads the aggregate CSV into memory on startup, and wires the scorer,
    chatbot, and HDX client together.
  \item[\texttt{scorer.py}] Core scoring engine.  All neglect-index formula
    logic (Equation~\ref{eq:neglect}) lives here.
  \item[\texttt{chatbot.py}] Claude API integration.  Runs an agentic loop
    where Claude can call tools (data queries, parameter updates, HDX
    lookups) before returning a natural language response.
  \item[\texttt{hdx\_client.py}] HTTP client for
    the HDX HAPI.  Provides eight data-fetching methods.
\end{description}

\subsection{Frontend (\texttt{frontend/src/})}

\begin{description}[labelindent=0.5cm, labelsep=1em]
  \item[\texttt{App.js}] Root component; holds all shared state
    (\code{csvData}, \code{currentParams}, chat history, selected crisis).
  \item[\texttt{CsvViewer.js}] Sortable, filterable data table.
    Supports cluster-level and country-level rollup views.
  \item[\texttt{TsneChart.js}] Interactive 2D scatter plot using Recharts.
    Points are coloured by neglect index; KMeans cluster assignments are
    overlaid; outliers are marked with red rings.
  \item[\texttt{Chat.js}] Chatbot panel.  Renders Claude's Markdown responses
    and automatically applies parameter updates returned by the API.
  \item[\texttt{Weights.js}] Collapsible methodology panel showing active
    weights as normalised bar charts.
  \item[\texttt{CounterfactualSlider.js}] ``What-if'' overlay.  A funding
    slider triggers a \apiendpoint{/counterfactual} query; the resulting rank
    change and new t-SNE position are reflected in real time.
  \item[\texttt{WorldMap.js}] Geographic view using \texttt{react-simple-maps};
    hover-to-filter integration with the table.
  \item[\texttt{RadarModal.js}] Detailed metrics popover for a selected crisis.
\end{description}

% -----------------------------------------------------------------------------
\section{API Reference}
% -----------------------------------------------------------------------------

All endpoints return JSON.  The base URL is \code{http://localhost:8000}.

\subsection{\apiendpoint{GET /metadata}}

Returns the set of available countries, clusters, and years in the loaded dataset.

\textbf{Response:}
\begin{lstlisting}[language={}]
{ "countries": [...], "clusters": [...], "years": [...] }
\end{lstlisting}

\subsection{\apiendpoint{GET /ranking}}

The primary ranking endpoint.  Accepts 25+ query parameters.

\subsubsection*{Filter parameters}

\begin{center}
\begin{tabular}{lll}
  \toprule
  Parameter & Type & Description \\
  \midrule
  \param{last\_years} & int & Only include data from the most recent $N$ years \\
  \param{cluster[]} & string[] & Restrict to specific humanitarian clusters \\
  \param{country[]} & string[] & Restrict to specific country codes (ISO3) \\
  \param{min\_people\_in\_need} & float & Minimum people-in-need threshold \\
  \param{max\_coverage} & float & Maximum funding coverage (0--1) \\
  \param{min\_neglect\_index} & float & Minimum neglect index filter \\
  \param{neglect\_type[]} & string[] & Filter by structural neglect type \\
  \param{min\_consecutive\_years} & int & Minimum consecutive underfunded years \\
  \bottomrule
\end{tabular}
\end{center}

\subsubsection*{Weight parameters}

\begin{center}
\begin{tabular}{lll}
  \toprule
  Parameter & Default & Description \\
  \midrule
  \param{severity\_weight} & 0.6 & Weight of severity vs.\ funding gap \\
  \param{funding\_gap\_weight} & 0.4 & Weight of funding gap \\
  \param{need\_weight} & 0.5 & Sub-weight: people in need \\
  \param{ipc\_weight} & 0.4 & Sub-weight: IPC food security \\
  \param{events\_weight} & 0.1 & Sub-weight: conflict events \\
  \bottomrule
\end{tabular}
\end{center}

\subsubsection*{Display and pagination parameters}

\begin{center}
\begin{tabular}{lll}
  \toprule
  Parameter & Default & Description \\
  \midrule
  \param{critical\_threshold} & 0.8 & Neglect index for ``critical'' label \\
  \param{high\_threshold} & 0.6 & Neglect index for ``high'' label \\
  \param{sort\_by} & \texttt{neglect\_index} & Column to sort by \\
  \param{sort\_desc} & true & Descending sort \\
  \param{limit} & 50 & Results per page \\
  \param{offset} & 0 & Pagination offset \\
  \param{n\_bootstrap} & 0 & Bootstrap iterations (0 = off, max 500) \\
  \bottomrule
\end{tabular}
\end{center}

\subsection{\apiendpoint{GET /tsne}}

Returns a 2D t-SNE projection of the scored crises, plus KMeans cluster assignments
(2--6 clusters) and outlier flags.

\subsection{\apiendpoint{GET /crisis/\{country\}/\{year\}/\{cluster\}}}

Returns full scoring context for a single crisis: all metrics, ranks, severity case
(A/B/C/D), and bootstrap confidence bounds.

\subsection{\apiendpoint{GET /counterfactual}}

Computes the hypothetical effect of additional funding on a crisis.

\textbf{Query parameters:} \param{country\_code}, \param{cluster},
\param{additional\_funding} (USD).

\textbf{Response:}
\begin{lstlisting}[language={}]
{
  "current_neglect_index": 0.84,
  "new_neglect_index":     0.71,
  "rank_change":           +12,
  "coverage_improvement":  0.18
}
\end{lstlisting}

\subsection{\apiendpoint{POST /chat} (SIDE TASK)}

Runs the Claude-powered analysis pipeline.

\textbf{Request body:}
\begin{lstlisting}[language={}]
{
  "messages": [ {"role": "user", "content": "..."} ],
  "current_params": { ...active frontend weights/filters... }
}
\end{lstlisting}

Claude may call up to ten internal tools before responding:
\begin{itemize}[nosep]
  \item \code{query\_ranking} --- execute a ranking query with custom params
  \item \code{update\_ranking\_parameters} --- push new params to the frontend
  \item \code{hdx\_search\_locations}, \code{hdx\_get\_humanitarian\_needs},
        \code{hdx\_get\_affected\_populations}, \code{hdx\_get\_food\_security},
        \code{hdx\_get\_conflict\_events}, \code{hdx\_get\_funding},
        \code{hdx\_get\_operational\_presence}
\end{itemize}

\textbf{Response:}
\begin{lstlisting}[language={}]
{
  "reply":             "Prose analysis...",
  "parameter_update":  { ...merged params dict or null... },
  "ranking_snapshot":  [ ...top results at time of response... ]
}
\end{lstlisting}

% -----------------------------------------------------------------------------
\section{Data Sources}
% -----------------------------------------------------------------------------

All primary data files reside in the \code{data/} directory.

\begin{longtable}{p{2cm}p{8.5cm}p{6cm}}
  \toprule
  \textbf{Source} & \textbf{File(s)} & \textbf{Description} \\
  \midrule
  \endhead
  \bottomrule
  \endfoot
  FTS (OCHA) &
  \texttt{fts\_requirements\_funding\_global.csv},
  \texttt{fts\_requirements\_\allowbreak funding\_globalcluster\_global} &
  Global funding and requirements, 2000--2025, by country, cluster, and year.
  Primary source for coverage calculations. \\[4pt]

  HPC / HNO &
  \texttt{hpc\_hno\_2024.csv},
  \texttt{hpc\_hno\_2025.csv} &
  Humanitarian Needs Overview: people in need (PiN) by country, cluster, and year.
  Source for \code{need\_rank}. \\[4pt]

  IPC (FEWS Net) &
  \texttt{ipc\_global\_area\_wide.csv} &
  Integrated Food Security Phase Classification by country and year.
  Population shares in phases 1--5.
  Used for \code{ipc\_severity} in the ``Food Security'' cluster. \\[4pt]

  ACLED &
  \texttt{number\_of\_events\_targeting\_civilians\_\ldots.csv} &
  Conflict events targeting civilians, by country and year.
  Source for \code{events\_rank}. \\[4pt]
\end{longtable}

\subsubsection*{Live data (HDX HAPI)}

When the environment variable \code{HDX\_APP\_IDENTIFIER} is set, the chatbot
can query the HDX Humanitarian API (\texttt{api.hpc.tools}) for real-time data
on needs, food security, conflict, funding, and operational presence.

\subsubsection*{Data merging pipeline}

The aggregate data frame construction follows these steps:

\begin{enumerate}[nosep]
  \item Load FTS plan-level and cluster-level funding.
  \item Load HNO people-in-need figures (2024, 2025).
  \item Merge FTS + HNO on \texttt{(countryCode, cluster, year)}.
  \item Left-merge IPC scores on \texttt{(countryCode, year)} where
        \texttt{cluster = "Food Security"}.
  \item Left-merge ACLED civilian events on \texttt{(countryCode, year)}.
  \item Left-merge structural neglect signals (computed from 6-year FTS window).
\end{enumerate}

% -----------------------------------------------------------------------------
\section{Installation \& Setup}
% -----------------------------------------------------------------------------

\subsection{Prerequisites}

\begin{itemize}[nosep]
  \item Python 3.11+
  \item Node.js 18+ and \code{npm}
  \item \code{uv} (Python package manager) or \code{pip}
  \item \code{latexmk} and a \LaTeX{} distribution (for this document)
\end{itemize}

\subsection{Backend}

\begin{lstlisting}[language=bash]
uv run uvicorn api.main:app
\end{lstlisting}

The server will be available at \url{http://localhost:8000}.
Interactive API docs (Swagger) are at \url{http://localhost:8000/docs}.

\subsection{Frontend}

\begin{lstlisting}[language=bash]
cd frontend
npm install
npm start   # -> http://localhost:3000
\end{lstlisting}

\subsection{Data Exploration}

The scripts used for data discovery and preliminary analyses was left in `scripts` directory.

\subsection{Environment Variables}

\begin{center}
\begin{tabular}{lll}
  \toprule
  Variable & Required & Description \\
  \midrule
  \code{ANTHROPIC\_API\_KEY} & Yes & Anthropic API key for the chatbot \\
  \code{HDX\_APP\_IDENTIFIER} & No & Base64-encoded \texttt{app\_name:app\_email}
                                     for live HDX data \\
  \bottomrule
\end{tabular}
\end{center}

% -----------------------------------------------------------------------------
\section{Usage Guide}
% -----------------------------------------------------------------------------

\subsection{Ranking Crises}

After starting both services, open \url{http://localhost:3000}.
The table on the bottom left displays all crises ranked by the neglect index using the
default weights ($S_w = 0.6$, $G_w = 0.4$).

\begin{itemize}[nosep]
  \item Use the \textbf{filter bar} to restrict by text, cluster, or neglect type.
  \item Toggle between \textbf{cluster view} (one row per country--cluster pair)
        and \textbf{country view} (rolled up to country level).
  \item Priority labels (\textcolor{critical-red}{critical} /
        \textcolor{high-orange}{high} / medium / low) appear in the first column.
\end{itemize}

\subsection{Adjusting Weights}

Open the \textbf{Weights} panel to see the active weight configuration.
Ask Claude in the chat panel to adjust weights (e.g., ``put more emphasis on
food security and reduce the conflict weight''); Claude will push updated
parameters to the frontend automatically.

\subsection{t-SNE Visualisation}

Switch to the \textbf{t-SNE view} to see all crises projected onto a 2D plane
where proximity reflects similarity of the scoring features.
Points are coloured from green (low neglect) to red (high neglect).
Outlier crises are ringed in red.

\subsection{Counterfactual Analysis}

Click any row in the table to open the \textbf{Counterfactual Slider}.
Drag the funding slider to see in real time:

\begin{itemize}[nosep]
  \item The new neglect index and rank after the hypothetical funding is added.
  \item The new position on the t-SNE chart (where would it land in the embedding space).
\end{itemize}

The aid can be utilized for \textbf{projecting the possible improvement} caused by an increased funding.

\subsection{Chatbot}

The chat panel accepts natural language queries such as:

\begin{itemize}[nosep]
  \item \textit{``Which food security crises in Sub-Saharan Africa have been underfunded for more than three consecutive years?''}
  \item \textit{``Show me the top 10 most neglected health crises and explain what drives their scores.''}
  \item \textit{``What would happen if the Yemen Food Security cluster received an extra \$200 million?''}
\end{itemize}

Claude will query the ranking engine, call HDX tools as needed, and return a
prose response.  If it suggests parameter changes, the frontend updates automatically.

% -----------------------------------------------------------------------------
\section{Limitations and Caveats}
% -----------------------------------------------------------------------------

\begin{enumerate}
  \item \textbf{Data coverage is uneven.}
    IPC food-insecurity data exists only for the ``Food Security'' cluster.
    ACLED events are available at country level, not cluster level---all clusters
    in the same country receive the same events rank.
    Crises with missing requirements or needs figures are not scored.

  \item \textbf{Funding data reflects reported figures, not actual spending.}
    FTS coverage ratios measure pledged/tracked funding against stated requirements.
    Under-reported funding or inflated requirements both distort the neglect index.

  \item \textbf{Bootstrap intervals are approximate.}
    The bootstrap procedure resamples the full crisis pool, not individual data
    points within a crisis.  Narrow clusters with few observations will exhibit
    wide intervals.

  \item \textbf{The scoring is relative, not absolute.}
    A country that scores ``critical'' is critical \emph{relative to its cluster peers}.
    A poorly covered cluster with uniformly low coverage may have many critical entries
    even if absolute need is modest.

  \item \textbf{The tool is for decision support, not automated allocation.}
    Rankings should be interpreted alongside qualitative context.
    Do not use the neglect index as the sole basis for funding decisions.
\end{enumerate}

% -----------------------------------------------------------------------------
\section{Project Structure}
% -----------------------------------------------------------------------------

\begin{lstlisting}[language={},numbers=none]
datathon-un/
+-- api/
|   +-- main.py            # FastAPI application and endpoints
|   +-- scorer.py          # Neglect index scoring engine
|   +-- chatbot.py         # Claude API integration
|   +-- hdx_client.py      # HDX HAPI HTTP client
+-- frontend/
|   +-- src/
|   |   +-- App.js                   # Root component / state
|   |   +-- CsvViewer.js             # Data table
|   |   +-- TsneChart.js             # t-SNE scatter plot
|   |   +-- Chat.js                  # Chatbot panel
|   |   +-- Weights.js               # Methodology weights display
|   |   +-- CounterfactualSlider.js  # Funding scenario modal
|   |   +-- WorldMap.js              # Geographic map view
|   |   +-- RadarModal.js            # Detailed metrics popover
|   +-- package.json
+-- data/                  # CSV and XLSX data files (gitignored)
+-- scripts/
|   +-- data/
|       +-- create_big_csv.py  # Data aggregation script
+-- notebooks/             # Exploratory Jupyter notebooks
+-- docs/
|   +-- documentation.tex  # This document
|   +-- workshop_notes.md  # Datathon brief and notes
+-- pyproject.toml         # Python dependencies
+-- .env.example           # Environment variable template
\end{lstlisting}

% -----------------------------------------------------------------------------
\appendix
\section{Humanitarian Cluster Taxonomy}
% -----------------------------------------------------------------------------

The IASC cluster system organises humanitarian response into thematic areas.
The following clusters appear in the HCNI dataset:

\begin{center}
\begin{tabular}{ll}
  \toprule
  Cluster & Description \\
  \midrule
  Food Security & Food assistance and livelihoods \\
  Health & Primary and secondary healthcare \\
  Nutrition & Acute malnutrition treatment and prevention \\
  WASH & Water, sanitation, and hygiene \\
  Shelter & Emergency shelter and NFIs \\
  Protection & Safety, dignity, and rights \\
  Education & Emergency and transitional education \\
  Logistics & Humanitarian supply chain \\
  Emergency Telecoms & Connectivity and ICT \\
  Camp Coordination & Site management for displacement \\
  Early Recovery & Restoration of basic services \\
  Multi-Sector & Refugee response (UNHCR-led) \\
  \bottomrule
\end{tabular}
\end{center}

\end{document}
